\section{The witness term and an overview of the proof}\label{section: overview}

We define
\begin{equation}\label{equation: Psi}
  \begin{aligned}
    \Kt    &\defeq \vb \inter \bigl(((\vr \compo \vr)^{*} \inter \vc) \compo \Dt\bigr), &
    \Nt    &\defeq \vb \setminus \Kt,\\
    \Ht    &\defeq (\St^{2} \compo \Kt \compo \St^{2}) \union (\St \compo \Kt \compo \St)
              \union (\St^{2} \compo \Nt \compo \St) \union (\St \compo \Nt \compo \St^{2}), &
    \Psit  &\defeq \Ht^{*}.
  \end{aligned}
\end{equation}
The star $(\vr \compo \vr)^{*}$ lies inside $\Kt$, and $\Kt$ lies inside the star $\Ht^{*}$.
All abbreviations are star-free, so $\sh(\Psit) = 2$.
Hence \Cref{theorem: main} follows from the next theorem.

\begin{theorem}\label{theorem: lower bound}
  No term $E$ with $\sh(E) \le 1$ is equivalent to $\Psit$ on $\Cls{2}$.
  This holds even if only $\RSUB$ valuations are considered in which $\vr, \vc, \vb$ are irreflexive.
\end{theorem}

The roles of the parts of $\Psit$ are as follows.
\begin{itemize}
  \item $\Kt$ reads one bit from an interval $\range{u, v}$ with $(u, v) \in \vb$, called a \emph{module}.
        The bit says whether some $\vc$-successor $y$ of $u$ can be reached from $u$ by an $\vr$-path of even length (\Cref{section: modules}).
  \item $\Ht$ shifts this bit to two ports before $u$ and two ports after $v$.
        A module in $\Kt$ swaps the ports and a module in $\Nt$ keeps them, so $\Ht^{*}$ computes a parity (\Cref{section: cells}).
\end{itemize}
Let $E$ be a term with $\sh(E) \le 1$.
Let $q$ be an upper bound of $\cd(F)$ for all subterms $F^{*}$ of $E$.
The proof of \Cref{theorem: lower bound} has three steps.
\begin{enumerate}
  \item\label{step: local pair}
        We construct two modules $\M^{+}_{q}, \M^{-}_{q}$ on the same vertices (\Cref{proposition: local pair}).
        They differ in a single $\vc$-pair, $\M^{+}_q$ is in $\Kt$ and $\M^{-}_q$ is not.
        However, no observation of the form $((A \inter \Dt) \compo F^{*} \compo (C \inter \Dt))(u, v)$ with $A, F, C \in \SF{q}$ distinguishes them.
        Note that $\Kt(u, v)$ is not of this form, since it intersects $(\vr \compo \vr)^{*}$ with $\vc$ at the pair $(u, y)$.
  \item\label{step: replacement}
        We put modules into a sequence of cells.
        A replacement lemma shows the following (\Cref{lemma: replacement}).
        For every star $F^{*}$ of $E$, replacing a module that lies strictly between the endpoints by $\M^{-}_{q}$ does not change the value of $F^{*}$.
  \item\label{step: logic}
        Hence the value of $E$ can be computed from the bit sequence of the cells by a first-order formula with modular predicates.
        That logic cannot express the parity computed by $\Psit$ (\Cref{section: lower bound}).
\end{enumerate}
Step~\ref{step: local pair} is the technical core.
The two modules must agree on all observations of depth $q$ at once.
An observation may read the prefix of a module up to some point and, at the same time, the suffix after that point.
Prefixes are described by first-order types, which form an aperiodic monoid.
Suffixes are described by a finite automaton, whose transition monoid may contain groups.
For example, $(\vr \compo \vr)^{*}$ counts modulo two.
We use an idempotent from the minimal ideal of the product monoid.
It acts as a reset on the first component and as a local identity on the second.
